\chapter{Gries-Schneider Quantification}

Here we gather key definitions from \cite{gries.93} 
with regard to boolean quantification.

First, Leibniz for quantification:

$\begin{array}{l@{\quad}l}
   (8.12)~\textbf{Leibniz:}
   &
   \inferrule
     {P = Q}
     {(\star x \mid E[x:=P] : S) = (\star x \mid E[z:=Q] : S)}
\\
\\ &
   \inferrule
     {R \implies (P=Q)}
     {(\star x \mid R : E[x:=P]) = (\star x \mid R : E[z:=Q])}
\end{array}$\newline\newline

The following are all \textbf{Axiom}s in \cite{gries.93}.

$\begin{array}{l@{\quad}l}
   (8.13)~\textbf{Empty Range:}
   &
   ( \star x \mid false : P) = u
\end{array}$\newline\newline

$\begin{array}{l@{\quad}l}
   (8.14)~\textbf{One-point rule:} \text{ Provided } \lnot occurs(`x',`E'),
\\ \\ \qquad 
   (\star x\mid x= E : P) = P[x:=E] \quad 
\end{array}$\newline\newline

$\begin{array}{l@{\quad}l}
   (8.15)~\textbf{Distributivity:} 
   \text{ Provided }P,Q:\Bool \text{ or } R \text{ is finite,}
\\ \\ \qquad
   (\star x \mid R : P) \star (\star x \mid R :Q) 
   = (\star x \mid R : P \star Q)
\end{array}$\newline\newline

$\begin{array}{l@{\quad}l}
   (8.16)~\textbf{Range split:}
   \text{ Provided } R \land S \equiv false
     \text{ and } P:\Bool \text{ or } R \text{ and } S \text{ are finite,}
\\ \\ \qquad
   (\star x \mid R \lor S : P)
   = (\star x \mid R : P) \star (\star x \mid S : P)
\end{array}$\newline\newline

$\begin{array}{l@{\quad}l}
   (8.17)~\textbf{Range split:}
   \text{ Provided } P:\Bool \text{ or } 
   R \text{ and } S \text{ are finite,}
\\ \\ \qquad
   (\star x \mid R \lor S : P) \star (\star x \mid R \land S : P)
   = (\star x \mid R : P) \star (\star x \mid S : P)
\end{array}$\newline\newline

$\begin{array}{l@{\quad}l}
   (8.18)~\textbf{Range split for idempotent $\star$ :}
 \\ \\ \qquad
   (\star x \mid R \lor S : P)
   = (\star x \mid R : P) \star (\star x \mid S : P)
\end{array}$\newline\newline

$\begin{array}{l@{\quad}l}
   (8.19)~\textbf{Interchange of dummies:}
   \text{ Provided } \star \text{ is idempotent or } \\
   R \text{ and } Q \text{ are finite, }  
   \lnot occurs(`y',`R'), \lnot occurs(`x',`Q')
\\ \\ \qquad
   (\star x \mid R : (\star y \mid Q : P)) 
   = (\star y \mid Q : (\star x \mid R : P))
\end{array}$\newline\newline

$\begin{array}{l@{\quad}l}
   (8.20)~\textbf{Nesting:} \text{ Provided } \lnot occurs(`y',`R'),
\\ \\ \qquad
   (\star x,y \mid R \land Q : P)
   = (\star x \mid R : (\star y \mid Q : P))
\end{array}$\newline\newline

$\begin{array}{l@{\quad}l}
   (8.21)~\textbf{Dummy renaming:} \text{ Provided } \lnot occurs(`y',`R,P'),
\\ \\ \qquad
   (\star x \mid R : P) = (\star y \mid R[x:=y] : P[x:=y])
\end{array}$\newline\newline

We note the following theorem:

$\begin{array}{l@{\quad}l}
   (8.23)~\textbf{Theorem Split off term:} \text{ For } n:\Nat
\\ \\ \qquad
   (\star i \mid 0 \leq i < n+1 : P)
   = (\star i \mid 0 \leq i < n : P) \star P[i:=n]
\\ \qquad
   (\star i \mid 0 \leq i < n+1 : P)
   = P[i:=0] \star (\star i \mid 0 < i < n+1 : P)
\end{array}$\newline\newline



More \textbf{Axiom}s.

~(9.1) $\qquad$ We re-write $(\land x \mid R : P)$ as $(\forall x \mid R : P)$.
\newline

$\begin{array}{l@{\quad}l}
   (9.2)~\textbf{Trading:} 
\\ \\ \qquad
   (\forall x \mid R : P) \equiv (\forall x\mid: R \implies P)
\end{array}$\newline\newline

$\begin{array}{l@{\quad}l}
   (9.5)~\textbf{Distributivity of $\lor$ over $\forall$:} 
   \text{ Provided }  \lnot occurs(`x',`P')
\\ \\ \qquad
   P \lor (\forall x \mid R : Q)
   = (\forall x \mid R: P \lor Q)
\end{array}$\newline\newline

$\begin{array}{l@{\quad}l}
   (9.17)~\textbf{Generalised De Morgan:} 
   \text{ Rewriting $(\lor x \mid R : P)$ as $(\exists x \mid R : P)$}
\\ \\ \qquad
   (\exists x \mid R : P) \equiv \lnot (\forall x \mid R : \lnot P)
\end{array}$\newline\newline

We have the following law that is not described as either an axiom or theorem:
No proof is supplied, or even hinted at.

$\begin{array}{l@{\quad}l}
   (9.13)~\textbf{Instantiation } 
   (\forall x \mid: P) \implies P[x:=E]
\end{array}$\newline\newline

Can we prove this?
\begin{eqnarray*}
   \lefteqn{(\forall x \mid: P) \implies P[x:=E]}
\\&=& \text{adding explicit range predicate}
\\&& (\forall x \mid true : P) \implies P[x:=E]
\\&=&  true = P \lor \lnot P, \text{choose } P \equiv x=E
\\&& (\forall x \mid x=E \lor x \neq E : P) \implies P[x:=E]
\\&=& \text{Idempotent range split}
\\&& (\forall x \mid x=E : P) \land (\forall x \mid x \neq E : P) 
     \implies P[x:=E]
\\&=&  (A\land B \implies C) \equiv (A\implies C)\lor(B\implies C)
\\&& ((\forall x \mid x=E : P) \implies P[x:=E])
     \lor
     ((\forall x \mid x \neq E : P) \implies P[x:=E])
\end{eqnarray*}
